\documentclass{article}
\usepackage{graphicx} % Required for inserting images
\usepackage{caption} % Required for \captionof
\usepackage{amsmath}
\usepackage{float}
\usepackage{comment}
\usepackage[T1]{fontenc}
\usepackage[utf8]{inputenc}
\usepackage{lmodern}
\usepackage{microtype}
\usepackage{xcolor}
\usepackage{framed}
\renewenvironment{leftbar}{%
  \def\FrameCommand{\textcolor{orange}{\vrule width 3pt} \hspace{10pt}}%
  \MakeFramed {\advance\hsize-\width \FrameRestore}}%
 {\endMakeFramed}

\title{Hőtani mérés}
\author{Kern Luca, Csongor Vincze}
\date{2026 Április 14}

\begin{document}

\maketitle

\section*{5. Feladat: Határozza meg egy inhomogén tömegeloszlású lemezből készült minta tehetetlenségi nyomatékát a súlypontján átmenő és a lemez síkjára merőleges tengelyre vonatkozóan!}

\[
T'^2_{\max} = \frac{4\pi^2}{k^*} \left[ \theta + \theta_x + m(r_0 + r_1)^2 \right],
\]

\[
T'^2_{\min} = \frac{4\pi^2}{k^*} \left[ \theta + \theta_x + m(r_0 - r_1)^2 \right],
\]

%fotó, minta!!!!

$m_{minta}= \text{kg}$
$r_{1}= \text{m}$
$r_{2}= \text{m}$

\begin{table}[H]
    \centering
    \begin{tabular}{|c|c|}
        \hline
        Szögelfordulás [°] & Lengésiidő [s] \\ \hline
         &  \\ \hline 
         &  \\ \hline
         &  \\ \hline 
         &  \\ \hline
         &  \\ \hline 
         &  \\ \hline
         &  \\ \hline
         &  \\ \hline
        
    \end{tabular}
    \caption{Lengési idők az első lyukkal}
    \label{tab:lengesiido_1}
\end{table}

\begin{table}[H]
    \centering
    \begin{tabular}{|c|c|}
        \hline
        Szögelfordulás [°] & Lengésiidő [s] \\ \hline
         &  \\ \hline 
         &  \\ \hline
         &  \\ \hline 
         &  \\ \hline
         &  \\ \hline 
         &  \\ \hline
         &  \\ \hline
         &  \\ \hline
        
    \end{tabular}
    \caption{Lengési idők a második lyukkal}
    \label{tab:lengesiido_2}
\end{table}

$\Theta = $


\end{document}